% =============================================================
%  CHAPTER: Theoretical Framework
%  Thesis: Hallucination Detection in Large Language Models
%  Status: Revised draft — publication quality
% =============================================================

This chapter connects the internal computation of decoder-only Transformer language models with an energy-based view of associative retrieval. The central motivation is that the hidden state entering the feed-forward sublayer can be treated as a query over structured parametric memory. This interpretation is not adopted as a strict architectural identity, but rather as a theoretically motivated modeling perspective supported by prior work on feed-forward layers and associative memory~\cite{geva2021transformer, ramsauer2020hopfield}. On that basis, the chapter introduces a Modern Hopfield formulation used to define retrieval distributions and divergence-based indicators for hallucination detection.

% =============================================================

\section{The Transformer Architecture}
\label{sec:transformer}
% ================================================================

The Transformer, introduced by \citet{vaswani2017attention}, replaced the sequential
processing of recurrent networks with a fully parallel, attention-based computation.
It has since become the architectural substrate of virtually all modern language
models. Understanding its internal mechanics is essential to this work: the
hallucination detection framework developed in Chapter~\ref{chap:methods} operates
directly on the per-token hidden states produced at each layer of the forward pass.
Figure~\ref{fig:transformer_block} provides an overview of the complete block
structure; the subsections below describe each component in turn.

% --- Figure 1: full Pre-LN decoder block ----------------------------
\begin{figure}[H]
    \centering
    \includegraphics[width=0.95\linewidth]{figures/fig1_transformer_block.pdf}
    \caption{%
        A single decoder-only Transformer block using Pre-Layer Normalisation
        (RMSNorm) and a SwiGLU feed-forward sublayer, as employed by Qwen2.5 and
        Llama~3. Dashed lines denote residual (skip) connections. The hidden state
        $\mathbf{h}_t^{(l)}$ at the output of the second RMSNorm—labelled
        \emph{probe point}—is the representational signal analysed in
        this work.%
    }
    \label{fig:transformer_block}
\end{figure}
% --------------------------------------------------------------------

% ----------------------------------------------------------------
\subsection{From Recurrent Networks to the Attention Mechanism}
\label{subsec:rnn_to_attention}
% ----------------------------------------------------------------

Before the Transformer, sequence modelling relied on Recurrent Neural Networks
(RNNs). At each time step an RNN takes the current input token and a hidden state
carried over from the previous step, producing an updated state that summarises the
sequence history \citep{tunstall2022natural}. Because each step depends on the
previous one, computation cannot be parallelised across positions, which makes
training on long sequences slow.

In neural machine translation, an encoder RNN compressed the source sentence into a
single fixed-length vector, which a decoder RNN then consumed token by token. The
fundamental weakness of this design is that the decoder has access only to the
\emph{last} encoder hidden state: for long inputs, compressing an arbitrarily long
sequence into a fixed-dimensional vector causes significant information loss \citep{tunstall2022natural}.

The attention mechanism resolved this bottleneck by granting the decoder direct
access to \emph{all} encoder hidden states at every generation step, weighting each
by its relevance to the current decoding context \citep{tunstall2022natural}. The
Transformer generalises this idea by making attention the primary computational
mechanism of the entire model, discarding recurrence altogether.

% ----------------------------------------------------------------
\subsection{Tokens, Embeddings, and Positional Encoding}
\label{subsec:embeddings}
% ----------------------------------------------------------------

A Transformer operates on sequences of discrete units called \emph{tokens}, obtained
by segmenting raw text through a tokenizer. Let $x_1, x_2, \ldots, x_T \in
\mathcal{V}$ denote a sequence of $T$ tokens drawn from vocabulary $\mathcal{V}$.

\paragraph{Token embeddings.}
Each token $x_t$ is mapped to a dense vector via a learned embedding matrix
$\mathbf{E} \in \mathbb{R}^{|\mathcal{V}| \times D}$, where $D$ is the model
dimension. This converts discrete symbols into a continuous vector space amenable to
gradient-based optimisation. The resulting embeddings are context-free---the same
vector is assigned regardless of surrounding tokens---with context-sensitivity
introduced later through self-attention (Section~\ref{subsec:self_attention}).

\paragraph{Positional encoding.}
Because the Transformer processes all positions in parallel, explicit positional
information must be injected. The original formulation \citep{vaswani2017attention}
used deterministic sinusoidal functions:
\begin{equation}
    \mathbf{p}_t^{(2k)} = \sin\!\left(\frac{t}{10000^{2k/D}}\right), \qquad
    \mathbf{p}_t^{(2k+1)} = \cos\!\left(\frac{t}{10000^{2k/D}}\right),
    \label{eq:sinusoidal_pe}
\end{equation}
where $k \in \{0, \ldots, D/2-1\}$ indexes the embedding dimension. Modern LLMs such as Qwen2.5 \citep{qwen2025qwen25} and Llama~3 \citep{grattafiori2024llama3}, which are the primary model families examined in this work, instead apply \emph{Rotary Positional Encoding} (RoPE) \citep{su2024roformer}. RoPE encodes relative position by rotating query and key vectors prior to their inner product, which improves generalisation to long sequences.

The initial hidden state matrix for a sequence of $T$ tokens is:
\begin{equation}
    \mathbf{H}^{(0)} = \mathbf{E}[x_1, \ldots, x_T] + \mathbf{P},
    \label{eq:initial_hidden}
\end{equation}
where $\mathbf{P} \in \mathbb{R}^{T \times D}$ collects the positional encodings and
$\mathbf{H}^{(0)} \in \mathbb{R}^{T \times D}$ is the input to the first
Transformer block.

% ----------------------------------------------------------------
% ----------------------------------------------------------------
\subsection{Self-Attention}
\label{subsec:self_attention}
% ----------------------------------------------------------------

Self-attention is a mechanism that produces a contextualised representation of each
token by computing a weighted combination of all token representations in the
sequence, where the weights reflect pairwise relevance. To compute these weights,
each token simultaneously takes three roles: it generates a \emph{query} vector
representing its informational need, a \emph{key} vector representing its content as
a retrieval target, and a \emph{value} vector carrying the content it contributes to
other tokens. These three vectors are obtained via independent learned linear
projections of the hidden state matrix $\mathbf{H} \in \mathbb{R}^{T \times D}$:
\begin{equation}
    \mathbf{Q} = \mathbf{H}\mathbf{W}^Q, \qquad
    \mathbf{K} = \mathbf{H}\mathbf{W}^K, \qquad
    \mathbf{V} = \mathbf{H}\mathbf{W}^V,
    \label{eq:qkv_projections}
\end{equation}
where $\mathbf{W}^Q, \mathbf{W}^K \in \mathbb{R}^{D \times d_k}$ and
$\mathbf{W}^V \in \mathbb{R}^{D \times d_v}$ are learned weight matrices. The
relevance of each key to a given query is measured by their scaled inner product;
the resulting scores are normalised via softmax and used to compute a weighted sum
of the value vectors:
\begin{equation}
    \text{Self-Attention}(\mathbf{H}) =
    \text{softmax}\!\left(\frac{\mathbf{Q}\mathbf{K}^\top}{\sqrt{d_k}}\right)\mathbf{V}.
    \label{eq:self_attention}
\end{equation}
The $\sqrt{d_k}$ scaling factor prevents the inner products from growing large in
high-dimensional spaces, which would push the softmax into regions of near-zero
gradient. The output is a context-aware hidden state matrix in which each token's
representation is a convex combination of all value vectors, with coefficients
determined by its query-key compatibilities. Unlike recurrent networks, where context
is accumulated sequentially, self-attention integrates information from any position
in a single parallelisable operation \citep{vaswani2017attention}.

\paragraph{Causal masking.}
In autoregressive LLMs, the prediction of token $x_{t+1}$ must condition only on
$x_1, \ldots, x_t$. This constraint is enforced during the forward pass by setting
the pre-softmax logits to $-\infty$ for all positions $j > t$, driving the
corresponding attention weights to zero. The resulting \emph{causal} or
\emph{masked} self-attention ensures that no token attends to future positions. All
models examined in this work apply causal masking throughout every layer, as
illustrated in Figure~\ref{fig:transformer_block}.

% ----------------------------------------------------------------
\subsection{Multi-Head Attention}
\label{subsec:mha}
% ----------------------------------------------------------------

A single attention head produces one weighted aggregation of the value vectors, which
limits it to capturing a single pattern of pairwise relevance per layer. Natural
language, however, requires reasoning over several types of dependency simultaneously:
a token may need to resolve syntactic agreement, semantic relatedness, and coreference
within the same layer. \citet{vaswani2017attention} addressed this by running $h$
attention heads in parallel, each with its own learned projections operating in a
lower-dimensional subspace of size $d_k = D/h$. The outputs are concatenated and
projected back to the model dimension (Figure~\ref{fig:transformer_block}):
\begin{equation}
    \text{MHA}(\mathbf{H}) =
    \text{Concat}(\text{head}_1, \ldots, \text{head}_h)\,\mathbf{W}^O,
    \qquad
    \text{head}_i = \text{Attention}\!\left(
        \mathbf{H}\mathbf{W}^Q_i,\;
        \mathbf{H}\mathbf{W}^K_i,\;
        \mathbf{H}\mathbf{W}^V_i
    \right),
    \label{eq:mha}
\end{equation}
where $\mathbf{W}^O \in \mathbb{R}^{h d_v \times D}$ is the output projection. Each
head specialises in a different relational pattern, and their combination yields a
richer contextualised representation than any single head could produce alone.% ----------------------------------------------------------------
\subsection{The Feed-Forward Sublayer}
\label{subsec:ffn}
% ----------------------------------------------------------------

Following multi-head attention, each block applies a \emph{position-wise
feed-forward network} (FFN) to every token representation independently. Unlike
attention, which mixes information across positions, the FFN transforms each hidden
state $\mathbf{h}_t^{(l)}$ in isolation, applying the same learned function to every
position within a given layer. Its role is therefore complementary: whereas attention
determines \emph{which} information is collected, the FFN determines \emph{how} that
information is re-encoded at each position. This position-wise independence makes
each $\mathbf{h}_t^{(l)}$ a clean unit of analysis, a property that this work
exploits directly.

In the original Transformer \citep{vaswani2017attention}, the FFN is a two-layer
network with ReLU activation:
\begin{equation}
    \mathrm{FFN}\!\left(\mathbf{h}_t^{(l)}\right)
    =
    \max\!\left(\mathbf{0},\; \mathbf{h}_t^{(l)}\mathbf{W}_1 + \mathbf{b}_1\right)
    \mathbf{W}_2 + \mathbf{b}_2,
    \label{eq:ffn_original}
\end{equation}
where $\mathbf{W}_1 \in \mathbb{R}^{D \times D_{\mathrm{ff}}}$ and
$\mathbf{W}_2 \in \mathbb{R}^{D_{\mathrm{ff}} \times D}$, with intermediate
dimension $D_{\mathrm{ff}}$ typically set to $4D$.

\paragraph{SwiGLU.}
\citet{dauphin2017language} introduced \emph{Gated Linear Units} (GLU) in the
context of gated convolutional language models: an element-wise product between a
sigmoid-gated projection and a linear projection, giving the network an explicit
mechanism for selective information flow. \citet{shazeer2020glu} subsequently
adapted this gating principle to the Transformer FFN and proposed a family of
variants by replacing the sigmoid gate with alternative activation functions. Among
these, SwiGLU has become the standard in modern LLMs—including Qwen2.5 and Llama~3,
as well as Mistral and Gemma—and is the variant relevant to this work:
\begin{equation}
    \mathrm{FFN}_{\mathrm{SwiGLU}}\!\left(\mathbf{h}_t^{(l)}\right)
    =
    \left(
        \mathrm{Swish}\!\left(\mathbf{h}_t^{(l)}\mathbf{W}_1\right)
        \odot
        \mathbf{h}_t^{(l)}\mathbf{V}
    \right)\mathbf{W}_2,
    \label{eq:swiglu}
\end{equation}
where $\mathbf{W}_1, \mathbf{V} \in \mathbb{R}^{D \times D_{\mathrm{ff}}}$ are the
gate and value projection matrices of the FFN (the symbol $\mathbf{V}$ here follows
\citealt{shazeer2020glu} and refers to the FFN value projection, distinct from the
attention value matrix in Equation~\ref{eq:qkv_projections}),
$\mathbf{W}_2 \in \mathbb{R}^{D_{\mathrm{ff}} \times D}$ is the output projection,
and $\mathrm{Swish}(z) = z \cdot \sigma(z)$ \citep{ramachandran2017searching}.
Figure~\ref{fig:swiglu} details the internal computational graph.
Because SwiGLU uses three weight matrices rather than two, $D_{\mathrm{ff}}$ is
reduced by a factor of $2/3$ relative to the standard $4D$ setting to maintain
comparable parameter counts \citep{shazeer2020glu}.

% --- Figure 2: SwiGLU internal structure ----------------------------
\begin{figure}[H]
    \centering
    \includegraphics[width=0.85\linewidth]{figures/fig2_swiglu_structure.pdf}
    \caption{%
        Internal structure of the SwiGLU feed-forward sublayer
        (Equation~\ref{eq:swiglu}). The input hidden state $\mathbf{h}_t^{(l)} \in
        \mathbb{R}^D$ is projected along two parallel branches: a gate branch through
        $\mathbf{W}_1 \in \mathbb{R}^{D \times D_{\mathrm{ff}}}$ followed by the Swish
        activation $\mathrm{Swish}(z) = z \cdot \sigma(z)$, and a value branch through
        $\mathbf{V} \in \mathbb{R}^{D \times D_{\mathrm{ff}}}$, where $D$ is the model
        dimension and $D_{\mathrm{ff}}$ is the intermediate dimension. Their element-wise
        product is projected back to $\mathbb{R}^D$ by $\mathbf{W}_2 \in
        \mathbb{R}^{D_{\mathrm{ff}} \times D}$. The three-matrix parameterisation
        distinguishes SwiGLU from the original two-matrix FFN and maps directly onto the
        key-value memory decomposition introduced in Section~\ref{sec:ffn_memory}.%
    }
    \label{fig:swiglu}
\end{figure}
% --------------------------------------------------------------------

The connection between the FFN parameterisation and associative memory retrieval is
made precise in Section~\ref{sec:ffn_memory}, where the matrices $\mathbf{W}_1$ and
$\mathbf{V}$ are identified with a key bank and a value bank respectively.

% ----------------------------------------------------------------
\subsection{Residual Connections and Layer Normalisation}
\label{subsec:residual_layernorm}
% ----------------------------------------------------------------

Each Transformer block composes the MHA and FFN sublayers through two structural mechanisms: residual connections and layer normalisation. Together, these mechanisms form the \emph{residual stream}, which is the progressively refined representation of the input shown in Figure~\ref{fig:transformer_block}.

\paragraph{Residual connections.}
A residual (skip) connection adds the sublayer input directly to its output:
\begin{equation}
    \mathbf{y} = \mathbf{x} + f(\mathbf{x}),
    \label{eq:residual}
\end{equation}
where $f(\cdot)$ is the sublayer transformation. This provides a direct gradient
path during backpropagation, mitigating vanishing gradients, and encourages each
sublayer to learn an incremental refinement of its input rather than a complete
transformation from scratch.

\paragraph{Pre-Layer Normalisation.}
Modern LLMs apply normalisation \emph{before} each sublayer (Pre-LN), as shown in
Figure~\ref{fig:transformer_block}. Using $\tilde{\mathbf{H}}^{(l)}$ to denote the
intermediate hidden state after the attention residual, the full block update is:
\begin{align}
    \tilde{\mathbf{H}}^{(l)} &= \mathbf{H}^{(l-1)}
        + \mathrm{MHA}^{(l)}\!\left(\mathrm{RMSNorm}\!\left(\mathbf{H}^{(l-1)}\right)\right),
    \label{eq:residual_attn} \\
    \mathbf{H}^{(l)} &= \tilde{\mathbf{H}}^{(l)}
        + \mathrm{FFN}^{(l)}\!\left(\mathrm{RMSNorm}\!\left(\tilde{\mathbf{H}}^{(l)}\right)\right).
    \label{eq:residual_ffn}
\end{align}
The token-level hidden state $\mathbf{h}_t^{(l)} =
\mathrm{RMSNorm}\!\left(\tilde{\mathbf{h}}_t^{(l)}\right)$ entering the FFN is the primary signal used in the detection framework of this work. This representation corresponds to the \emph{probe point} shown in Figure~\ref{fig:transformer_block}.

\paragraph{RMSNorm.}
Both Qwen2.5 and Llama~3 employ Root Mean Square Normalisation (RMSNorm), a
simplified variant that omits the mean-centring step of standard layer normalisation:
\begin{equation}
    \mathrm{RMSNorm}\!\left(\mathbf{h}_t^{(l)}\right)
    =
    \frac{\mathbf{h}_t^{(l)}}
         {\sqrt{\frac{1}{D}\sum_{d=1}^{D}\!\left(h_{t,d}^{(l)}\right)^2 + \epsilon}}
    \odot \boldsymbol{\gamma},
    \label{eq:rmsnorm}
\end{equation}
where $\epsilon > 0$ ensures numerical stability and $\boldsymbol{\gamma} \in
\mathbb{R}^D$ is a learned scale vector. Removing mean-centring reduces
computational overhead while retaining the stabilising effect on activation
magnitudes.

% ----------------------------------------------------------------
\subsection{Taxonomy of Transformer Architectures}
\label{subsec:transformer_taxonomy}
% ----------------------------------------------------------------

Transformer-based models are commonly grouped into three families according to their
attention pattern \citep{tunstall2022natural}. \emph{Encoder-only} models process
the full input sequence bidirectionally and are primarily used for classification
and language understanding tasks. \emph{Encoder-decoder} models learn a
sequence-to-sequence mapping through a cross-attention bridge between a bidirectional
encoder and a causal decoder, suited to tasks such as machine translation. The
present work is concerned exclusively with \emph{decoder-only} models, which employ
causal masked self-attention to autoregressively generate one token at a time. This family includes Qwen2.5, Llama 3, Mistral, and Gemma, which are the architectures examined in the experiments of Chapter~\ref{chap:experiments}. Their autoregressive structure
produces a sequence of per-token hidden states $\mathbf{h}_t^{(l)}$ at each layer
$l$ and position $t$, which constitute the representational signal that the
detection mechanism analyses.
% =============================================================
% =============================================================
% ================================================================
\section{Internal Representations of LLMs}
\label{sec:internal_representations}
% ================================================================

The previous section established how a decoder-only Transformer builds token
representations layer by layer through self-attention and feed-forward sublayers
connected by residual additions. Using the notation established there,
$\mathbf{h}_t^{(l)} \in \mathbb{R}^D$ denotes the hidden state of token $t$ after
layer $l$, and $\mathbf{H}^{(l)} \in \mathbb{R}^{T \times D}$ the corresponding
matrix across all positions. Across layers and positions, these vectors define the
internal representational trajectory of the model during generation.

Not all of these states play the same role in this work. The objective here is not
to treat hidden representations as generic features, but to study them as queries
over parametric memory. This requires identifying the sub-computation in which a
token representation is matched against stored parameters and transformed according
to that match. In decoder-only Transformers, that role is played by the feed-forward
sublayer.

% ----------------------------------------------------------------
\subsection{Feed-Forward Layers as Parametric Memory}
\label{sec:ffn_memory}
% ----------------------------------------------------------------

Each feed-forward sublayer processes token representations independently, applying
the same learned transformation to every position. This position-wise independence
means that the hidden state $\mathbf{h}_t^{(l)}$ acts on the sublayer's parameters
without interacting with any other token, making it a natural unit for studying
how individual representations interact with learned knowledge.

\citet{geva2021transformer} formalised this observation by showing that the
feed-forward sublayer can be written as an unnormalised key-value memory lookup.
Letting $\mathbf{r} \in \mathbb{R}^D$ denote the token representation entering the
sublayer (corresponding to $\mathbf{h}_t^{(l)}$ in the notation of
Equation~\ref{eq:residual_ffn}), the computation takes the form:
\begin{equation}
    \mathrm{FFN}^{(l)}(\mathbf{r})
    =
    g\!\left(\mathbf{r}\,{\mathbf{K}^{(l)}}^{\!\top}\right)\mathbf{V}^{(l)},
    \label{eq:ffn_kv_view}
\end{equation}
where $\mathbf{K}^{(l)} \in \mathbb{R}^{K_l \times D}$ is the first weight matrix,
whose $K_l$ rows act as key vectors, $\mathbf{V}^{(l)} \in \mathbb{R}^{K_l \times D}$
is the second weight matrix, whose rows act as value vectors, and $g(\cdot)$ is an
element-wise nonlinearity such as ReLU. The inner products $\mathbf{r}{\mathbf{K}^{(l)}}^{\!\top}$
measure the alignment between the input and each key, the nonlinearity $g$ produces
a coefficient for each memory slot, and $\mathbf{V}^{(l)}$ assembles the output as
a weighted combination of value vectors. The distinction from a normalised memory is
that the canonical Transformer does not apply softmax to these coefficients, making
the retrieval unnormalised \citep{geva2021transformer}.

This framing has a concrete empirical basis. \citet{geva2021transformer} showed that
individual key vectors correlate with human-interpretable input patterns: lower
layers tend to capture shallow surface patterns such as recurring n-grams, while
upper layers encode more semantic content. The corresponding value vectors promote
output distributions concentrated on tokens that plausibly follow those patterns.
The feed-forward sublayer thus acts as a distributed, layer-specific pattern library,
and the input representation selectively activates entries in that library through
its inner products with the stored keys.

% ----------------------------------------------------------------
\subsection{Gated Feed-Forward Layers and the Memory Bank}
\label{subsec:gated_ffn_memory}
% ----------------------------------------------------------------

The models examined in this work — Qwen2.5 and Llama~3 — use the SwiGLU variant
described in Section~\ref{subsec:ffn} rather than the ReLU-based FFN of
Equation~\ref{eq:ffn_kv_view}. In generic form, a gated feed-forward layer is:
\begin{equation}
    \mathrm{FFN}^{(l)}(\mathbf{r})
    =
    \Bigl(
        \phi\!\left(\mathbf{r}\,{\mathbf{A}^{(l)}}^{\!\top}\right)
        \odot
        \mathbf{r}\,{\mathbf{B}^{(l)}}^{\!\top}
    \Bigr)\mathbf{C}^{(l)},
    \label{eq:gated_ffn_generic}
\end{equation}
where $\mathbf{A}^{(l)}, \mathbf{B}^{(l)} \in \mathbb{R}^{K_l \times D}$,
$\mathbf{C}^{(l)} \in \mathbb{R}^{K_l \times D}$, $\phi(\cdot)$ is a smooth
activation, and $\odot$ denotes element-wise multiplication. In the SwiGLU
parameterisation of Equation~\ref{eq:swiglu}, $\mathbf{A}^{(l)} = \mathbf{W}_1$,
$\mathbf{B}^{(l)} = \mathbf{V}$, and $\mathbf{C}^{(l)} = \mathbf{W}_2$ (see
Figure~\ref{fig:swiglu}).

The gating structure means that the effective retrieval score of each intermediate
unit depends on two projections of the input simultaneously, rather than on a single
inner product with a stored key. Any identification of a fixed memory bank with a
specific weight matrix is therefore a modelling choice rather than a structural
theorem. This is an important distinction: the key-value view of
\citet{geva2021transformer} applies directly to the ReLU case, and requires a
deliberate extension when applied to gated architectures.

In this thesis, the memory bank of layer $l$ is defined as the collection of row
vectors of $\mathbf{A}^{(l)}$:
\begin{equation}
    \mathbf{M}^{(l)} \coloneqq \mathbf{A}^{(l)}
    =
    \begin{bmatrix}
        {\mathbf{m}_1^{(l)}}^{\!\top} \\
        \vdots \\
        {\mathbf{m}_{K_l}^{(l)}}^{\!\top}
    \end{bmatrix}
    \in \mathbb{R}^{K_l \times D}.
    \label{eq:memory_bank_definition}
\end{equation}
This choice is motivated by the role of $\mathbf{A}^{(l)}$ in the gating
computation: the term $\phi(\mathbf{r}{\mathbf{A}^{(l)}}^{\!\top})$ determines
which intermediate units are selectively activated, making $\mathbf{A}^{(l)}$ the
matrix whose rows most directly function as pattern detectors in the sense of
\citet{geva2021transformer}. The matrices $\mathbf{B}^{(l)}$ and $\mathbf{C}^{(l)}$
remain essential to the full computation, but do not serve the same pattern-matching
role.

With this definition, the hidden state $\mathbf{r} = \mathbf{h}_t^{(l)}$ entering
the feed-forward sublayer is treated as a query over the memory bank
$\mathbf{M}^{(l)}$: the inner products $\mathbf{h}_t^{(l)}{\mathbf{M}^{(l)}}^{\!\top}$
measure the alignment of the token's representation with each stored pattern.
Section~\ref{sec:hopfield} formalises this retrieval operation within a Modern
Hopfield Network framework, which provides an energy function and a closed-form
retrieval distribution over the memory bank.
% =============================================================
% ================================================================
\section{Modern Hopfield Networks}
\label{sec:hopfield}
% ================================================================

Hopfield networks provide a principled framework for associative memory retrieval:
given a query state and a collection of stored patterns, they define an energy
function whose minima correspond to retrieved patterns, and an update rule that
descends that energy. Their relevance to this thesis is not historical. They supply
the mathematical language needed to convert the alignment between a hidden state and
a layer-specific memory bank into an energy, a retrieval distribution, and
concentration measures that can be linked to hallucination behaviour in
Section~\ref{sec:divergences}.

% ----------------------------------------------------------------
\subsection{From Classical to Modern Associative Retrieval}
\label{subsec:classical_to_modern_hopfield}
% ----------------------------------------------------------------

Classical Hopfield networks \citep{hopfield2007hopfield} store binary patterns in a
symmetric weight matrix $\mathbf{W} = \sum_{i=1}^{K} \mathbf{m}_i\mathbf{m}_i^\top
\in \mathbb{R}^{D \times D}$ constructed via Hebbian learning, and retrieve them by
iterating a sign update that decreases the quadratic energy:
\begin{equation}
    E_{\mathrm{cl}}(\mathbf{q})
    =
    -\tfrac{1}{2}\mathbf{q}^{\top}\mathbf{W}\mathbf{q},
    \qquad \mathbf{q} \in \{-1,+1\}^D.
    \label{eq:classical_hopfield_energy}
\end{equation}
This formulation is too restrictive for the present setting: states are binary, and
the storage capacity grows only linearly with the dimension $D$.

\citet{ramsauer2020hopfield} introduced a modern Hopfield network with continuous
states and exponential storage capacity. Let
\begin{equation}
    \mathbf{M}
    =
    \begin{bmatrix}
        \mathbf{m}_1^\top\\
        \vdots\\
        \mathbf{m}_K^\top
    \end{bmatrix}
    \in \mathbb{R}^{K \times D}
    \label{eq:generic_memory_bank}
\end{equation}
be a memory bank whose $K$ rows are stored patterns, and let $\mathbf{q} \in
\mathbb{R}^D$ be a continuous query state. The modern Hopfield energy is:
\begin{equation}
    E(\mathbf{q};\,\mathbf{M},\beta)
    =
    -\frac{1}{\beta}\log\sum_{i=1}^{K}
    \exp\!\bigl(\beta\,\mathbf{m}_i^{\top}\mathbf{q}\bigr)
    +\frac{1}{2}\|\mathbf{q}\|_2^{2}
    + C(\mathbf{M},\beta,K),
    \label{eq:modern_hopfield_energy}
\end{equation}
where $\beta > 0$ is an inverse temperature parameter and
$C(\mathbf{M},\beta,K) = \beta^{-1}\log K + \frac{1}{2}\max_i\|\mathbf{m}_i\|_2^2$
is an additive constant that depends only on the memory bank, not on the query. When
computing retrieval distributions or divergences between distributions, this constant
cancels and can be omitted.

The energy is minimised by the update rule \citep{ramsauer2020hopfield}:
\begin{equation}
    \mathbf{q}^{+}
    =
    \mathbf{M}^{\top}\,\mathrm{softmax}\!\bigl(\beta\,\mathbf{M}\mathbf{q}\bigr).
    \label{eq:modern_hopfield_update}
\end{equation}
The updated state is a convex combination of the stored patterns, with weights given
by a softmax over their inner products with the current query. \citet{ramsauer2020hopfield}
prove global convergence and identify three types of fixed points: a global fixed
point near the arithmetic mean of all patterns (reached when no pattern strongly
dominates, corresponding to a near-uniform softmax), metastable states averaging
over a correlated subset of patterns, and single-pattern fixed points reached when
the query is well-separated from all other stored patterns. This trichotomy will
be relevant when interpreting entropy as a retrieval-concentration signal in
Section~\ref{subsec:retrieval_distribution_llm}.

\paragraph{Connection to the attention mechanism.}
\citet{ramsauer2020hopfield} show that the update rule of
Equation~\eqref{eq:modern_hopfield_update} is equivalent to scaled dot-product
attention under the identification $\beta = 1/\sqrt{d_k}$, with the stored patterns
serving simultaneously as keys and values. Concretely, if the stored patterns
$\mathbf{m}_i$ play the role of both key and value vectors, then
Equation~\eqref{eq:modern_hopfield_update} reduces to:
\begin{equation}
    \mathbf{q}^{+}
    =
    \mathrm{Attention}\!\left(\mathbf{q},\,\mathbf{M},\,\mathbf{M}\right)
    =
    \mathbf{M}^{\top}\,\mathrm{softmax}\!\!\left(
        \frac{\mathbf{M}\mathbf{q}}{\sqrt{d_k}}
    \right).
    \label{eq:mhn_attention_equiv}
\end{equation}
This is the attention mechanism of \citet{vaswani2017attention} with $K = V =
\mathbf{M}$. It is not the same as self-attention, where queries, keys, and values
arise from separate learned projections of the same sequence. The connection is
nonetheless significant: it shows that the Transformer architecture and modern
Hopfield networks share the same core retrieval operation, which justifies using
Hopfield machinery to analyse the internal behaviour of Transformer layers.

\paragraph{Sparse extensions.}
Later work generalised the dense softmax retrieval map to sparse alternatives such
as sparsemax and $\alpha$-entmax, yielding sparse modern Hopfield networks with
stronger convergence to single-pattern fixed points \citep{martins2023sparse}. This
thesis retains the dense softmax formulation because it produces a full probability
distribution over all stored patterns, which is precisely the object needed for the
entropy and divergence-based analysis that follows.

% ----------------------------------------------------------------
\subsection{Retrieval Distribution Induced by an LLM Hidden State}
\label{subsec:retrieval_distribution_llm}
% ----------------------------------------------------------------

The quantity of interest in this thesis is not the converged fixed point of the
Hopfield dynamics. Rather, it is the retrieval distribution produced by a single
application of the softmax map to a hidden state, capturing the alignment structure
of that state with the memory bank at the moment of generation.

Instantiating the generic framework of Section~\ref{subsec:classical_to_modern_hopfield}
with the Transformer quantities established in Sections~\ref{sec:internal_representations}
and \ref{subsec:gated_ffn_memory}, the query is the hidden state entering the FFN
at layer $l$ and token position $t$:
\begin{equation}
    \mathbf{q} = \mathbf{h}_t^{(l)},
    \label{eq:query_identification}
\end{equation}
and the memory bank is the layer-specific matrix $\mathbf{M}^{(l)}$ defined in
Equation~\eqref{eq:memory_bank_definition}. The induced retrieval distribution is:
\begin{equation}
    p_i^{(l,t)}
    =
    \frac{
        \exp\!\bigl(\beta\,{\mathbf{m}_i^{(l)}}^{\!\top}\mathbf{h}_t^{(l)}\bigr)
    }{
        \sum_{j=1}^{K_l}
        \exp\!\bigl(\beta\,{\mathbf{m}_j^{(l)}}^{\!\top}\mathbf{h}_t^{(l)}\bigr)
    },
    \qquad i = 1, \ldots, K_l.
    \label{eq:retrieval_distribution}
\end{equation}
This distribution converts the hidden state into a probability vector over the $K_l$
stored patterns of layer $l$. The retrieved prototype — the pattern that the hidden
state most nearly recovers — is the expectation under this distribution:
\begin{equation}
    \bar{\mathbf{m}}_t^{(l)}
    =
    \sum_{i=1}^{K_l} p_i^{(l,t)}\,\mathbf{m}_i^{(l)}.
    \label{eq:retrieved_pattern}
\end{equation}

Equation~\eqref{eq:retrieval_distribution} is the central bridge between the
Transformer and the Hopfield formalism: it turns a token-level hidden state into an
interpretable distribution over learned memory patterns. The shape of that
distribution carries information about how the model is using its parametric
knowledge at a given layer and position.

By the fixed-point analysis of \citet{ramsauer2020hopfield}, a near-uniform
distribution corresponds to the global averaging regime, where no single pattern is
strongly preferred. A distribution concentrated on a few patterns corresponds to a
metastable state. A distribution sharply peaked on one pattern indicates that the
query is well within the basin of attraction of that memory. A natural summary of
this concentration is the retrieval entropy:
\begin{equation}
    H^{(l,t)}
    =
    -\sum_{i=1}^{K_l} p_i^{(l,t)} \log p_i^{(l,t)}.
    \label{eq:retrieval_entropy}
\end{equation}
High entropy indicates diffuse, unstructured retrieval; low entropy indicates that
the hidden state is sharply aligned with a small subset of stored patterns. The
hypothesis motivating this work is that correct and hallucinated generation differ
systematically in the entropy and divergence structure of the distributions
$\bigl\{p_i^{(l,t)}\bigr\}$ across layers. Section~\ref{sec:divergences}
introduces the information-theoretic tools used to quantify that divergence.

% =============================================================
% ================================================================
\section{Information-Theoretic Divergences}
\label{sec:divergences}
% ================================================================

The retrieval distribution $\{p_i^{(l,t)}\}$ of Equation~\eqref{eq:retrieval_distribution}
and the entropy $H^{(l,t)}$ of Equation~\eqref{eq:retrieval_entropy} characterise
the alignment structure of a single hidden state with the memory bank at one layer
and one position. The hallucination detection hypothesis of this work requires
comparing two such distributions: one induced during the question prefill phase and
one induced during answer generation. This section introduces the four
information-theoretic quantities used to quantify that dissimilarity. Throughout,
$P = (p_1, \ldots, p_K)$ and $Q = (q_1, \ldots, q_K)$ denote two probability
distributions over $K$ outcomes; these are distribution labels and are distinct from
the query and key matrices of Section~\ref{subsec:self_attention}.

% ----------------------------------------------------------------
\subsection{Shannon Entropy}
\label{subsec:shannon_entropy}
% ----------------------------------------------------------------

The Shannon entropy of a discrete distribution $P$ \citep{shannon1948mathematical}
measures the average uncertainty of drawing from it:
\begin{equation}
    H(P) = -\sum_{i=1}^{K} p_i \log p_i,
    \label{eq:shannon_entropy}
\end{equation}
It attains its maximum $\log K$ at the uniform
distribution and its minimum $0$ at a point mass. In the present context,
$H^{(l,t)}$ in Equation~\eqref{eq:retrieval_entropy} is the instance of
Equation~\eqref{eq:shannon_entropy} with $P = \{p_i^{(l,t)}\}_{i=1}^{K_l}$,
measuring how spread the retrieval is across the $K_l$ patterns of layer $l$.

Because $K_l$ varies across architectures and configurations, the analysis uses the
\emph{normalised entropy}:
\begin{equation}
    \bar{H}(P) = \frac{H(P)}{\log K},
    \label{eq:normalised_entropy}
\end{equation}
which maps the range of $H$ to $[0,1]$ and makes per-layer and cross-model
comparisons dimensionless. A value of $\bar{H} \approx 1$ indicates near-uniform
retrieval; $\bar{H} \approx 0$ indicates near-deterministic retrieval concentrated
on a single pattern.

% ----------------------------------------------------------------
\subsection{Kullback-Leibler Divergence}
\label{subsec:kl_divergence}
% ----------------------------------------------------------------

The Kullback-Leibler (KL) divergence from $Q$ to $P$
\citep{kullback1951information} measures the expected log-likelihood ratio of $P$
over $Q$ under $P$:
\begin{equation}
    D_{\mathrm{KL}}(P \| Q) = \sum_{i=1}^{K} p_i \log \frac{p_i}{q_i}.
    \label{eq:kl_divergence}
\end{equation}
By Gibbs' inequality, $D_{\mathrm{KL}}(P \| Q) \geq 0$, with equality if and only
if $P = Q$ almost everywhere. The KL divergence is asymmetric: $D_{\mathrm{KL}}(P
\| Q) \neq D_{\mathrm{KL}}(Q \| P)$ in general. In the application of this work,
both distributions are softmax outputs over the full memory bank, ensuring strictly
positive entries and a well-defined KL divergence. The asymmetric form is retained
here as the building block for the symmetric measures introduced below; the
direction of the asymmetry carries meaning when one distribution serves as a fixed
reference and the other as a test.

% ----------------------------------------------------------------
\subsection{Jensen-Shannon Divergence}
\label{subsec:js_divergence}
% ----------------------------------------------------------------

The Jensen-Shannon divergence (JSD) \citep{lin1991divergence} symmetrises the KL
divergence via the mixture $M = \tfrac{1}{2}(P + Q)$:
\begin{equation}
    D_{\mathrm{JS}}(P \| Q)
    =
    \frac{1}{2}\,D_{\mathrm{KL}}(P \| M)
    +
    \frac{1}{2}\,D_{\mathrm{KL}}(Q \| M).
    \label{eq:js_divergence}
\end{equation}
An equivalent expression in terms of entropies makes the connection to
Equation~\eqref{eq:shannon_entropy} explicit:
\begin{equation}
    D_{\mathrm{JS}}(P \| Q)
    =
    H(M) - \frac{1}{2}H(P) - \frac{1}{2}H(Q).
    \label{eq:js_entropy_form}
\end{equation}
The JSD is symmetric, always finite, and bounded: $0 \leq D_{\mathrm{JS}}(P \| Q)
\leq \log 2$ with natural logarithms. These properties make it the primary
divergence used in this work for comparing retrieval distributions across the
prefill and generation phases.

% ----------------------------------------------------------------
\subsection{Hellinger Distance}
\label{subsec:hellinger}
% ----------------------------------------------------------------

The Hellinger distance provides a geometric complement to
the entropy-based JSD. It is defined as:
\begin{equation}
    H^2(P, Q)
    =
    \frac{1}{2}\sum_{i=1}^{K}
    \!\left(\sqrt{p_i} - \sqrt{q_i}\right)^{\!2},
    \label{eq:hellinger_squared}
\end{equation}
where $H^2(P, Q)$ is the squared Hellinger distance, also called the Hellinger
divergence. It satisfies $0 \leq H^2(P, Q) \leq 1$, with $H^2 = 0$ iff $P = Q$
and $H^2 = 1$ iff $P$ and $Q$ have disjoint support. The Hellinger distance is
symmetric, bounded regardless of whether $P$ and $Q$ share support, and can be
computed stably by operating on square-root probabilities rather than log-ratios,
avoiding the numerical instabilities that arise near zero-probability events.

The relationship between the Hellinger distance and the Bhattacharyya coefficient
$\mathrm{BC}(P, Q) = \sum_{i=1}^{K} \sqrt{p_i q_i}$ is direct:
\begin{equation}
    H^2(P, Q) = 1 - \mathrm{BC}(P, Q),
    \label{eq:hellinger_bhattacharyya}
\end{equation}
which shows that the Hellinger distance measures how far the two distributions are
from having identical square-root profiles. Used alongside JSD, it provides a
second, geometrically distinct measure of the shift between retrieval distributions
across generation phases, with complementary sensitivity to the tails of the
distribution.

In this work, both $D_{\mathrm{JS}}$ and $H^2$ are computed per generated token per
layer between the generation-phase retrieval distribution and the prefill reference
distribution. The specific definition of the prefill reference and the aggregation
of these per-token, per-layer scalars into sample-level features is formalised in
Chapter~\ref{chap:methods}.% =============================================================
% =============================================================
%  END OF CHAPTER
% =============================================================